% sec-03-ind.tex - Section 3, IND.  A main section.  DRAFT stage.
% Four figure slots, Figures 6 to 9, and three tables, Tables 8, 9 and 10.
% Every drafting instruction names trial-ind/final-ind/publication as the source
% stage 7 must read, and quotes the exact section file to quote from.
\section{IND}
\label{sec:ind}

\subsection{What was filed, and under what authority}

The Investigational New Drug application deposited July 1, 2026 supports a Phase
1, first-in-human, combined IND and Investigational Device Exemption
investigation of on-premises large language model directed eight-arm robotic
pancreaticoduodenectomy with perioperative daraxonrasib (RMC-6236) in patients
with KRAS G12-mutated pancreatic ductal adenocarcinoma \cite{phase1ind}. It is
filed under 21 CFR \S 312 for the investigational drug and coupled to a
significant-risk early-feasibility device study under 21 CFR \S 812 for the
LLM-directed robotic platform.

\draftinstr{Stage 7: quote directly from
\appfile{trial-ind/final-ind/publication/LaTeX Source Files.zip},
sections/sec-02-introduction.tex, the sentence beginning ``This Investigational
New Drug (IND) application'' and the sentence stating that the two pathways are
coordinated within a single integrated protocol because the device conducts the
resection that generates the operative and pathology inputs to the daraxonrasib
restart advisory. Attribute the quotation to the deposited IND, not to this
paper.}

\figslot{diagrams (python)-type / clustered infrastructure}{9.0}{Figure 6. Four
source clusters feeding twelve IND modules, with the three modules\\that had no
prior-system counterpart and were composed rather than copied.}

Figure 6 shows what the dossier was assembled from. Four clusters of existing
material fed it: adapted regulatory text, two clinical protocols with two
published outcome sets, a simulation record, and a device specification. Three
modules had no prior-system counterpart to copy, and were composed: the Physical AI
Subpart J text, the autonomy disclosure, and the verification cost ledger.

\draftinstr{Stage 7: build Table 8 from the twelve section files of
\appfile{trial-ind/final-ind/publication}. Columns: Module, ReGARDD position,
Character count at deposit, Figures carried. The IND carries 22 grayscale
figures in a single document sequence and section-scoped table numbering; state
both.}

\begin{apptable}
\begin{tabularx}{\textwidth}{>{\raggedright\arraybackslash}p{5.2cm} >{\raggedright\arraybackslash}p{2.2cm} >{\raggedright\arraybackslash}p{2.2cm} >{\raggedright\arraybackslash}X}
\headrow \thead{Module} & \thead{Position} & \thead{Characters} & \thead{What it establishes}\\
\midrule
Cover letter & 1 & to be filled & Transmittal and sponsor identity \\
Table of contents & 2 & to be filled & ReGARDD content order \\
FDA forms 1571 and 1572 & 3 & to be filled & The codified form set \\
Introduction & 4 & to be filled & 312.23(a)(3)(i) introductory statement \\
General investigational plan & 5 & to be filled & 312.23(a)(3)(iv) plan \\
Investigator's brochure & 6 & to be filled & 312.23(a)(5) brochure \\
Proposed clinical research & 7 & to be filled & 312.23(a)(6) protocols \\
Chemistry, manufacturing, control & 8 & to be filled & 312.23(a)(7) CMC \\
Pharmacology and toxicology & 9 & to be filled & 312.23(a)(8) nonclinical \\
Previous human experience & 10 & to be filled & 312.23(a)(9) prior exposure \\
Additional information & 11 & to be filled & 312.23(a)(10) \\
Relevant information & 12 & to be filled & 312.23(a)(11) \\
\end{tabularx}
\end{apptable}
\vspace{-0.6cm}
\tabcap{Table 8. The twelve modules of the deposited IND in ReGARDD content\\order,
with the codified requirement each is written to satisfy.}

\subsection{The clock}

\figslot{mermaid-type / gantt}{8.2}{Figure 7. Twelve IND modules on one clock:
the new system in hours across four\\days, and the prior system's published
calendar for the same work, in months.}

Figure 7 places both regimes on one axis, and Table 9 gives the day each
module closed against the prior system's name for the same work.

\draftinstr{Stage 7: expand this subsection using the four-stage build record in
\appfile{trial-ind/final-ind/README.md}: mermaid catalog, draft, full, final.
State the deposit date, July 1, 2026, and the repository version at deposit,
v4.3.0. Build Table 9 with columns Module, Repository source read, Generation
window, Prior-system counterpart.}

\begin{apptable}
\begin{tabularx}{\textwidth}{>{\raggedright\arraybackslash}p{4.0cm} >{\raggedright\arraybackslash}p{4.4cm} >{\raggedright\arraybackslash}p{2.0cm} >{\raggedright\arraybackslash}X}
\headrow \thead{Module} & \thead{Repository source read} & \thead{Window} & \thead{Prior-system counterpart}\\
\midrule
Introduction and plan & Adapted 21 CFR 312, ICH E6(R3) & Day 1 & Regulatory drafting, months \\
Investigator's brochure & Phase 1 protocol, RASolute 302 & Day 2 & Nonclinical write-up \\
Proposed clinical research & Phase 1 and Phase 2 protocols & Day 2 & Regulatory drafting \\
Chemistry, manufacturing, control & Public daraxonrasib record & Day 2 & CMC compilation \\
Pharmacology and toxicology & QSP and digital twin record & Day 3 & Nonclinical write-up \\
Previous human experience & RASolute 302, Dutch cohort & Day 3 & Internal QC \\
Figures, tables, back matter & The 22-figure catalog & Day 4 & Sponsor sign-off \\
\end{tabularx}
\end{apptable}
\vspace{-0.6cm}
\tabcap{Table 9. Each IND module, the repository source the generation read, the\\day
it closed, and the prior system's name for the same work.}

\subsection{Complexity and volume}

\figslot{d2-type / sql tables}{7.4}{Figure 8. Twelve codified IND content items
joined to the section file that\\satisfies each, and to the repository path
where that file is deposited.}

Figure 8 is the crosswalk a reviewer needs: each codified content item, the
generated file that satisfies it, and the deposited path. Table 10 states the
volume the crosswalk covers.

\draftinstr{Stage 7: build Table 10 from the compiled IND. Include total
character count across the twelve section files, figure count, table count,
reference count, the number of distinct 21 CFR cross-references, and the number
of distinct cited works. Compare each against a prior-system Phase 1 IND
prepared by a group.}

\begin{apptable}
\begin{tabularx}{\textwidth}{>{\raggedright\arraybackslash}p{5.0cm} >{\raggedright\arraybackslash}p{3.0cm} >{\raggedright\arraybackslash}X}
\headrow \thead{Volume measure} & \thead{Deposited IND} & \thead{What it means for a reviewer}\\
\midrule
Section files & 12 & One per ReGARDD content position \\
Figures, single document sequence & 22 & Every figure numbered once, end to end \\
Tables, section-scoped numbering & to be filled & Table 3.1, Table 5.2, and so on \\
Total characters across sections & to be filled & The dossier's written weight \\
Distinct 21 CFR cross-references & to be filled & Codified coverage \\
Authors & 1 & With model assistance, disclosed \\
Elapsed production time & 4 days & Against months in the prior system \\
\end{tabularx}
\end{apptable}
\vspace{-0.6cm}
\tabcap{Table 10. The volume one author reached in four days, measured on seven\\axes,
and what each measure means to a regulatory reviewer.}

\figslot{graphviz-type / fault tree}{8.6}{Figure 9. The clinical hold as a top
event over eight basal failures, with the\\two AND gates that make a hold require
a coincidence rather than a fault.}

Figure 9 states the sponsor's risk as logic rather than as a list. The two
clinical branches are AND gates, so a hold on clinical grounds requires both a
signal and a failure to respond to it. The documentation branch is an OR gate,
where any one omission suffices; that is precisely the branch a system which
regenerates a complete dossier in hours is built to eliminate.

\draftinstr{Stage 7: close the section by connecting the OR branch of Figure 9
to the four-day regeneration claim, and cite \appfile{phase1ind} and
\appfile{phase1guidance}. This is a main section: target the same character
count as sections 4, 5, 6 and 7.}
